% 19 November 2026. Drafted for Edward Baker with OpenAI Codex assistance.
% AGC-E02; dispatched configuration gpt-6.1-sol, effort ultra.
% Dispatch metadata is recorded, not independently attested backend metadata.
\section{A counted pair and its place in the gap chain}
\label{sec:pair-chain}

We derive the local laws directly from the physical equation
$H_t=-H_{zz}$, with the normalization of Section~\ref{sec:source}.
The derivation also specifies which information is lost when a pair is
viewed in isolation. These are inherited structural identities, rederived
from \cite{VCDCMManuscript,VCDCMDiscussion08}; their validity does not
constitute a new estimate preventing a collision.

\subsection{Complete products and finite parts}

The complete-divisor statements below use the following properties on a
compact real-time interval $J$. The field is real and even in $z$,
$H(t,0)\ne0$, and its even order-at-most-one product is
\begin{equation}
 H(t,z)=H(t,0)\prod_{\{\rho,-\rho\}}
                  \left(1-\frac{z^2}{\rho(t)^2}\right).
 \label{agc:pair-product}
\end{equation}
Every zero occurrence, including multiplicity, is retained; one factor
represents two opposite occurrences. The canonical growth estimates give,
uniformly on $J$,
\begin{equation}
 n_t(R)=O_J(R\log(2+R)),\qquad
 \sum_{|\rho|>R}|\rho|^{-2}=O_J(\log(2+R)/R).
 \label{agc:pair-count}
\end{equation}
The second assertion follows from the first by partial summation.
It gives locally uniform product convergence, and Cauchy's formula gives
convergence of every fixed spatial derivative. Evenness removes the linear
exponential in the order-one Hadamard product. This restriction matters:
the divisor alone does not determine logarithmic jets of an arbitrary
entire function with an unrestricted zero-free exponential factor.

Write $B=H_z/H$. At a simple zero $x(t)$, implicit differentiation gives
\begin{equation}
 x'=\frac{H_{zz}(t,x)}{H_z(t,x)}=2U_x,
 \qquad
 U_x=\lim_{z\to x}\left(B(t,z)-\frac1{z-x}\right).
 \label{agc:pair-node}
\end{equation}
For a finite polynomial the same equation reads
$x'=2\sum_{\rho\ne x}(x-\rho)^{-1}$, by differentiating its
factorization. For the complete field the finite part in
\eqref{agc:pair-node}, or the product grouping in
\eqref{agc:pair-product}, is the definition. Separate unregularized
first-power sums need not converge absolutely.

\subsection{Symmetric coordinates through a double zero}

Suppose a fixed positively oriented contour $\Gamma$ is zero-free and
counts exactly two zeros at a reference time. On a sufficiently small
time chart this remains so. With
\[
 J_k(t)=\frac1{2\pi\ii}\oint_\Gamma z^k B(t,z)\,\dd z,
 \qquad c=\frac{J_1}{2},\qquad D=\frac{J_2}{2}-c^2,
\]
the counted quadratic and local factorization are
\begin{equation}
 Q_*(t,z)=(z-c)^2-D,\qquad H(t,z)=Q_*(t,z)A(t,z),
 \qquad R(t,z)=\frac{A_z}{A}.
 \label{agc:pair-deflation}
\end{equation}
Here $A$ is jointly holomorphic and nonzero inside the contour. The
argument-principle formulas show that $c,D$ are holomorphic even when
the counted zeros merge. For a conjugation-invariant pair at real time,
$c,D$ are real. Write $z_\pm=c\pm s$, $s^2=D$, on a separated local
chart. Then $D>0$ describes two real roots, $D<0$ a conjugate pair,
and $D=0$ a double root. On a real adjacent edge the gap is
$g=z_+-z_-=2s>0$, so $D=g^2/4$.

\begin{proposition}[Pair law with the complete exterior]
\label{agc:pair-law-prop}
Let $R[a,b]=(R(b)-R(a))/(b-a)$, continued as $R_z(a)$ at
$a=b$, and set $M=-R[c-s,c+s]$. Then
\begin{equation}
 c'=R(t,c-s)+R(t,c+s),\qquad
 D'=2-4DM.
 \label{agc:pair-law}
\end{equation}
Under \eqref{agc:pair-product}--\eqref{agc:pair-count},
\begin{equation}
 M=\sum_{\rho\ \mathrm{exterior}}
       \frac1{(\rho-c)^2-D}.
 \label{agc:pair-M}
\end{equation}
This sum is absolute, includes every exterior occurrence, and continues
through exterior multiplicities and a double target while the contour
remains cleared. At a target double, $M=-R_z(t,c)$ and $D'=2$.
\end{proposition}

\begin{proof}
At a simple target, differentiating $H=Q_*A$ twice in space gives
\[
 z_+'=\frac1s+2R(t,c+s),\qquad
 z_-'=-\frac1s+2R(t,c-s).
\]
Taking the mean gives $c'$. Since $D=s^2$,
$D'=s(z_+'-z_-')=2+2s(R(c+s)-R(c-s))$, which is
\eqref{agc:pair-law}. These symmetric expressions extend holomorphically
to $s=0$ by the counted coordinates.
For an exterior occurrence, its logarithmic pole contributes
\[
 -\frac1{2s}\left(\frac1{c+s-\rho}-\frac1{c-s-\rho}\right)
       =\frac1{(\rho-c)^2-D}.
\]
The product regularizers are constant in space and cancel in this
difference. The resulting summand is $O(|\rho|^{-2})$ at infinity,
so \eqref{agc:pair-count} justifies summation and passage to the limit.
Alternatively, for $Q_*=(z-c)^2-D$ the same expression is
\begin{equation}
 M=-\frac1{2\pi\ii}\oint_\Gamma Q_*^{-1}B\,\dd z.
 \label{agc:pair-M-contour}
\end{equation}
For each exterior pole the integral is minus its rational weight;
the target contribution has total residue zero because all its poles
are inside $\Gamma$ and it decays as $z^{-3}$ at infinity.
The fixed contour proves holomorphy at all allowed multiplicities.
\end{proof}

The analytic exterior factor is invisible to polynomial-weighted count
integrals: $\oint_\Gamma z^kR\,\dd z=0$ for $k\ge0$.
Nevertheless its values at the targets enter their velocities through
the product rule. Identical current target positions can therefore have
different subsequent motion when their complete exteriors differ.
This is a local statement about deflation, not permission to change
the canonical field by an arbitrary analytic multiplier.

The first spatial variation of this analytic field can be retained
without discarding its effect on the clock. Put $Z=z-c$ and
$b=R_z(t,c)$. On a cleared contour enclosing the pair and its center,
Cauchy's formula and
$1/(Z^2-D)=1/Z^2+D/[Z^2(Z^2-D)]$ give
\begin{equation}
 \begin{split}
 D'&=2+4bD+E_D,\\
 b&=\frac1{2\pi\ii}\oint_\Gamma\frac{B}{Z^2}\,\dd z,\qquad
 E_D=\frac{4D^2}{2\pi\ii}\oint_\Gamma
                         \frac{B}{Z^2(Z^2-D)}\,\dd z.
 \end{split}
 \label{agc:pair-affine}
\end{equation}
Replacing $R$ by $B$ is valid since the added target rational functions
have all poles inside and total residue zero. If $|R_{zzz}|\le C_3$
on both center-to-root segments, Taylor's formula with integral
remainder gives $|E_D|\le(2/3)C_3|D|^2$.
The constant exterior response affects $c'$, while the affine slope
$b$ affects $D'$ already at order $D$. A small $E_D$ alone supplies
no bound on the uncertainty in $b$; the full rational kernel is exactly
$4D/Z^2+4D^2/[Z^2(Z^2-D)]=4D/(Z^2-D)$.

\subsection{Cumulative distances and the dipole response}

Now assume the \emph{complete} divisor is real and simple on an interval,
and label its ordered roots $x_n$, $n\in\mathbb Z$. Define
$g_n=x_{n+1}-x_n>0$, $c_n=(x_n+x_{n+1})/2$,
$D_n=g_n^2/4$, and let $M_n$ be \eqref{agc:pair-M} for this edge.
For $m\ge1$, put
\[
 L_{n,m}=\sum_{j=1}^m g_{n-j},\qquad
 R_{n,m}=\sum_{j=1}^m g_{n+j}.
\]
The left exterior root $x_{n-m}$ is at distances $-L_{n,m}$ and
$-(L_{n,m}+g_n)$ from the two targets. The right root
$x_{n+1+m}$ is at distances $R_{n,m}+g_n$ and $R_{n,m}$.
They exhaust the exterior, so absolute convergence proves
\begin{equation}
 \begin{split}
 M_n&=\sum_{m\ge1}\left\{
 \frac1{L_{n,m}(L_{n,m}+g_n)}+
 \frac1{R_{n,m}(R_{n,m}+g_n)}\right\},\\
 g_n'&=\frac4{g_n}-2g_nM_n.
 \end{split}
 \label{agc:pair-chain-law}
\end{equation}
This is a nonlocal equation: a gap enters many rows through cumulative
distances. It determines relative motion, leaving the common translation
coordinate to the node or center law. If only a finite real block is
known, only its actual terms have this positive-gap form. Every omitted
real or nonreal occurrence must remain in the complete interaction.

For targets $a,b=a+g$ and an exterior $\rho$,
\begin{equation}
 \frac1{(\rho-b)(\rho-a)}
   =\frac1g\left(\frac1{\rho-b}-\frac1{\rho-a}\right),\qquad
 M=-\frac{R(b)-R(a)}g=-\frac1g\int_a^b R_z(z)\,\dd z.
 \label{agc:pair-dipole}
\end{equation}
Thus the dipole interpretation is an exact difference of reciprocal
fields. At $w=\rho-(a+b)/2$ it is
$w^{-2}(1+O(g^2/w^2))$ as $|w|/g\to\infty$; the common far
response cancels before summation. $M$ is the complete response per
unit separation, not a population average.

There is also an exact spatial cancellation across edges. From
$x_n'=2U_n$ and \eqref{agc:pair-chain-law},
\begin{equation}
 \begin{split}
 g_nM_n&=\frac2{g_n}-(U_{n+1}-U_n),\\
 \sum_{n=a}^{b}g_nM_n
   &=2\sum_{n=a}^{b}\frac1{g_n}-(U_{b+1}-U_a).
 \end{split}
 \label{agc:pair-block}
\end{equation}
The proof is subtraction of the two node velocities followed by finite
telescoping. Shared interior force errors cancel if the same node
approximation is used on both edges. The identity controls a block
aggregate; it does not stop that aggregate from concentrating into a
single small gap. At a merger, the regular chart of the merging pair
may survive, but its neighboring edge charts lose exterior separation
because a formerly exterior occurrence meets a target.

\subsection{Changing density and local fluctuations}

The canonical chain has a varying logarithmic mean scale
\cite{VCDCMDiscussion09,VCDCMDiscussion10}. In the normalization
$x=2\gamma$ at $t=0$, the smooth counting main term and its derivative are
\begin{equation}
 \overline N(X)=\frac{X}{4\pi}\left(\log\frac{X}{4\pi}-1\right),
 \qquad \nu(X)=\frac1{4\pi}\log\frac{X}{4\pi},\qquad
 d_X=\nu(X)^{-1}.
 \label{agc:pair-density}
\end{equation}
Here $\nu$ differentiates the smooth main term, not the discrete
counting staircase or its error. The complete counting law alone
does not make the divisor real or give each real gap a lower bound.

Freeze a positive reference length $d$ near a chosen height and put
$\alpha_n=g_n/d$, $\mu_n=d^2M_n$, and
$\tau=(t-t_*)/d^2$. Rescaling the cumulative distances in
\eqref{agc:pair-chain-law} yields its exact dimensionless version, and
\begin{equation}
 D_n'=2-\alpha_n^2\mu_n,\qquad
 \frac{\dd\alpha_n}{\dd\tau}=\frac4{\alpha_n}-2\alpha_n\mu_n.
 \label{agc:pair-unfolded}
\end{equation}
Both the target fluctuation $\alpha_n$ and the environmental response
$\mu_n$ matter. A bound $d^2M_n\simeq2$ does not imply near cancellation
unless $g_n/d$ is controlled too. If $d$ instead varies with time,
$Y=D/d(t)^2$ obeys $Y'=D'/d^2-2Yd'/d$.
The logarithmic scale therefore cannot be inserted as a moving unit
without its derivative payment.

In the existing positive-time canonical acquisition, individual zeros
are first localized near a nonuniform clock; that input already gives
the target gap a positive lower bound. Only then are the nearby terms
compared with reference cumulative distances and the complete distant
tail paid to obtain $d^2M\simeq2$ \cite{VCDCMDiscussion10}.
The downstream interaction estimate is therefore not a derivation of
individual separation from smooth counting density alone. Extending
its domain requires additional canonical information.

As a normalization control only, a uniform infinite chain with gap $g$
has $M=2g^{-2}\sum_{m\ge1}[m(m+1)]^{-1}=2/g^2$ and $g'=0$.
Its telescoping sum checks the factors in the physical law; it is not
a canonical model of the height-dependent fluctuations.
Even the isolated quadratic heat field $z^2-2(t-T)$ obeys
\eqref{agc:pair-law}, with $M=0$, and has a double at $t=T$.
Such controls fail the literal theta-source premise, not this local
algebra. The next section organizes the complete exterior response;
the source estimate needed for continuation remains additional input.
